% Part: applied-modal-logics
% Chapter: epistemic-logic
% Section: truth-at-w

\documentclass[../../../include/open-logic-section]{subfiles}

\begin{document}

\olfileid[tr]{aml}{el}{trw}

\olsection{Bir Dünyada Doğruluk}

Normal modal mantıkta olduğu gibi her epistemik model de içindeki hangi
!!{formula}s{}in hangi dünyalarda doğru sayılacağını belirler. Bağıntısal
semantiğin temel kavramı için aynı ``$\mModel{M}$ modeli, !!{formula}~$!A$'yı
$w$ dünyasında doğru kılar'' gösterimini kullanırız. Bu bağıntı tümevarımlı
olarak tanımlanır ve modal olmayan bütün işleçler için normal modal mantık
durumuyla aynıdır.

\begin{defn}\ollabel{defn:mmodels}
  \emph{!!a{formula}~$!A$'nın $w$'deki doğruluğu}, bir $\mModel M$ içinde,
  yani $\mSat{M}{!A}[w]$ olmak üzere tümevarımlı olarak şöyle tanımlanır:
  \begin{enumerate}
  \tagitem{prvFalse}{\indcase{!A}{\lfalse}{$\mSat{M}{\lfalse}[w]$ hiçbir zaman geçerli değildir.}}{}
  \tagitem{prvTrue}{\indcase{!A}{\ltrue}{$\mSat{M}{\ltrue}[w]$ her zaman geçerlidir.}}{}
  \item $\mSat{M}{p}[w]$ ancak ve ancak $w \in V(p)$ ise geçerlidir.
  \tagitem{prvNot}{\indcase{!A}{\lnot !B}{$\mSat{M}{\indfrm}[w]$ ancak ve ancak
    $\mSat/{M}{!B}[w]$ ise geçerlidir.}}{}
  \tagitem{prvAnd}{\indcase{!A}{(!B \land !C)}{$\mSat{M}{\indfrm}[w]$ ancak ve ancak
    $\mSat{M}{!B}[w]$ ve $\mSat{M}{!C}[w]$ ise geçerlidir.}}{}
  \tagitem{prvOr}{\indcase{!A}{(!B \lor !C)}{$\mSat{M}{\indfrm}[w]$ ancak ve ancak
    $\mSat{M}{!B}[w]$ ya da $\mSat{M}{!C}[w]$ ise geçerlidir} (ikisi de olabilir).}{}
  \tagitem{prvIf}{\indcase{!A}{(!B \lif !C)}{$\mSat{M}{\indfrm}[w]$ ancak ve ancak
    $\mSat/{M}{!B}[w]$ ya da $\mSat{M}{!C}[w]$ ise geçerlidir.}}{}
  \tagitem{prvIff}{\indcase{!A}{(!B \liff !C)}{$\mSat{M}{\indfrm}[w]$ ancak ve ancak
    hem $\mSat{M}{!B}[w]$ hem $\mSat{M}{!C}[w]$ geçerliyse ya da
    ne $\mSat{M}{!B}[w]$ ne de $\mSat{M}{!C}[w]$ geçerliyse doğrudur.}}{}
  \item\ollabel{defn:sub:mmodels-box}
    \indcase{!A}{\Knows_a !B}{$\mSat{M}{\indfrm}[w]$ ancak ve ancak
    $\mSat{M}{!B}[w']$ bütün $w' \in W$ için $R_a ww'$ olduğunda
    geçerliyse doğrudur}
  \end{enumerate} 
\end{defn}

Ne var ki erişilebilirlik bağıntılarımıza getirilecek kısıtları tam burada
düşünmemiz gerekir. Nitekim \olref{defn:sub:mmodels-box} koşuluna göre,
!!a{formula}~$\Knows_a !B$, $w$'de, hiçbir $w'$ için $R_a ww'$ geçerli
olmadığında doğrudur. Bu, normal modal mantıktaki koşulun aynısıdır: bir dünyanın
ardılı yoksa bütün $\Box$-formülleri o dünyada boşsal olarak doğrudur. Fakat
$\Knows$ işlecinin \emph{bilgiyi} temsil ettiğini düşünürsek bu son derece
sezgiye aykırı görünür. Bir öznenin \emph{hiçbir} durumda hem $!A$'yı hem
$\lnot !A$'yı aynı anda bilemeyeceğini düşünmeye eğilimliyizdir.

Bir çözüm, epistemik mantıktaki erişilebilirlik bağıntımızın her zaman
\emph{yansımalı} olmasını güvenceye almaktır. Bu kabaca, edimsel dünyanın bir
öznenin bilgisiyle bağdaşması düşüncesine karşılık gelir. Epistemik mantıklar
aslında tipik olarak S5 kullanır; ancak tam olarak neyi temsil etmeleri
istendiğine bağlı olarak $\Knows_a$ bağıntısı için daha zayıf sistemler de
kullanılabilir.

\begin{figure}
  \begin{center}
    \begin{tikzpicture}[modal]
      \node[world] (w1) [label={[align=right]right:\mTrue{p}\\ \mFalse{q}}]{$w_1$}; 
      \node[world] (w2) [label={[align=right]right:\mTrue{p}\\ \mTrue{q}},
        above right=of w1]{$w_2$}; 
      \node[world] (w3) [label={[align=right]right:\mFalse{p}\\ \mFalse{q}},
        below right=of w1] {$w_3$};
      \draw[<->] (w1) to node {$a$} (w2);
      \draw[->,loop left] (w1) to node {$a,b$} (w1);
      \draw[<->] (w1) to [swap] node {$b$} (w3);
      \draw[->,loop above] (w2) to node {$a,b$} (w2) ;
      \draw[->,loop below] (w3) to node {$a,b$} (w3) ;
    \end{tikzpicture}
  \end{center}
  \caption{Basit bir epistemik model.}
  \ollabel{fig:simple}
\end{figure}

\begin{prob}
  Aşağıdakilerden hangilerinin \olref[aml][el][trw]{fig:simple} içinde
  geçerli olduğunu belirleyin:
  \begin{enumerate}
    \item $\mSat{M}{\lnot q}[w_1]$;
    \item $\mSat{M}{\Knows_a \lnot q}[w_1]$;
    \item $\mSat{M}{\Knows_b \lnot q}[w_1]$;
    \item $\mSat{M}{\Knows_b q \lor \Knows_b \lnot q}[w_2]$;
    \item $\mSat{M}{\Knows_a( \Knows_b q \lor \Knows_b \lnot q)}[w_2]$;
    \item $\mSat{M}{\EKnows_{\{a,b\}} \lnot q}[w_3]$;
    \end{enumerate}
\end{prob}

Bir dünyada doğruluğun temel tanımını verdiğimize göre modal geçerlilik ve
gerektirme gibi normal modal mantığın diğer semantik kavramları da modal
işleçlerin bu yeni yorumlanma biçimine uygulanarak doğrudan aktarılır.

Artık ortak bilgi işleci $\CKnows_G$ için de doğruluk koşulları verebiliriz.
\olref[sfr][rel][ops]{sec}'den, \emph{geçişli kapanış}~$R^+$'ın bir $R$
bağıntısı için şöyle tanımlandığını hatırlayın:
\begin{align*}
  R^+ &= \bigcup_{ n \in \mathbb{N}} R^n, 
\intertext{burada}
  R^0 & = R \text{ ve}\\ 
  R^{n+1} & = \Setabs{\tuple{x, z}}{\exists y (R^n xy \land Ryz)}.
\end{align*}
O hâlde $G$ bir özne grubu olduğunda, bütün öznelerin erişilebilirlik
bağıntılarının birleşiminin geçişli kapanışı olmak üzere $R_G = ( \bigcup_{b
\in G} R_b )^+$ tanımını veririz.

\begin{defn}
$G' \subseteq G$ ise $\mSat{M}{\CKnows_{G'} !A}[ w]$ ancak ve ancak her
$w'$ için, $R_{G'} w w'$ olduğunda $\mSat{M}{!A}[w']$ ise geçerlidir.
\end{defn}

\end{document}
